\begin{statementbox}
\par\smallskip\begin{center}\begin{tikzpicture}[x=1cm,y=0.86cm,>=Stealth,font=\small]
\path[use as bounding box] (-2,-3.5) rectangle (10,4.2);
\node[anchor=west,font=\small\bfseries] at (-1.8,3.8) {两切线交于 $P$，红线连接两个切点};
\draw[->,gray] (-1.7,0)--(3.3,0) node[right] {$x$};
\draw[->,gray] (0,-3.1)--(0,3.1) node[above] {$y$};
\node[below left] at (0,0) {$O$};
\draw[blue,thick,domain=-3:3,samples=81] plot ({\x*\x/4},\x);
\draw[orange!85!black,thick] (-1.4,-.4)--(1.8,2.8);
\draw[orange!85!black,thick] (-1.4,.4)--(1.8,-2.8);
\draw[red!75!black,very thick] (1,-2.7)--(1,2.7);
\fill (-1,0) circle (2pt) node[above left,yshift=12pt] {$P(-1,0)$};
\fill (1,2) circle (2pt) node[below right,xshift=5pt,yshift=-3pt] {$A(1,2)$};
\fill (1,-2) circle (2pt) node[above right,xshift=5pt,yshift=3pt] {$B(1,-2)$};

\node[anchor=north west,align=left,text width=5.7cm] at (4,2.7) {橙线：$PA,PB$ 是两条切线。\\[6pt]红线：$AB:x=1$ 是切点弦所在直线。\\[6pt]$P\notin AB$；切点弦连接的是 $A,B$。};
\node[anchor=west,blue] at (-1.6,-3.25) {$y^2=4x$（精确示例，$p=2$）};
\end{tikzpicture}\end{center}\par\smallskip

\textbf{\textcolor{CStatement}{一、已知条件与基本结论}}
\begin{description}[style=nextline, leftmargin=2.8em, labelsep=0.5em, itemsep=0.4em]
  \item[\textbf{条件：}] 抛物线 $C:y^2=2px\ (p>0)$，已知点 $P(x_0,y_0)$。从 $P$ 向抛物线作两条不同的实切线，切点为 $A,B$；直线 $AB$ 称为\textbf{切点弦所在直线}。
  \item[\textbf{两切点存在的充要条件：}] $y_0^2-2px_0>0$。
  \item[\textbf{结论一（$y^2$ 型）：}]
  \[
     \boxed{\ y_0y=p(x+x_0)\ }.
  \]
  \item[\textbf{结论二（$x^2$ 型）：}] 对 $C:x^2=2py\ (p>0)$，若 $x_0^2-2py_0>0$，则
  \[
     \boxed{\ x_0x=p(y+y_0)\ }.
  \]
\end{description}

\medskip
\textbf{\textcolor{CStatement}{二、使用公式之前，先判断有没有两条实切线}}
\[
\begin{array}{c|c|c}
\text{判断值} & \text{实切线条数} & \text{公式的实际含义}\\ \hline
>0 & 2 & \text{两个不同切点，存在切点弦}\\
=0 & 1 & \text{两个切点重合，所得直线是该点切线}\\
<0 & 0 & \text{无实切点，所得直线不是切点弦}
\end{array}
\]
对 $y^2=2px$，判断值取 $y_0^2-2px_0$；对 $x^2=2py$，判断值取 $x_0^2-2py_0$。

\medskip
\textbf{\textcolor{CStatement}{三、便于理解的附加结论}}
对 $y^2=2px$，当两切点存在且其中点为 $M$ 时，\textbf{$M$ 与 $P$ 的纵坐标相等}：
\[
  \boxed{\ y_M=y_0,\qquad x_M=\frac{y_0^2}{p}-x_0\ }.
\]
对于 $x^2=2py$，相应地有 $x_M=x_0$、$y_M=x_0^2/p-y_0$。

\medskip
\textbf{\textcolor{CStatement}{四、从直线的“一个点和斜率”理解公式}}
设 $A(x_1,y_1),B(x_2,y_2)$，$M(x_m,y_m)$ 为弦中点。两条切线都经过 $P$，分别列式后：
\[
\begin{aligned}
 \text{两式相减：}\quad& y_0(y_1-y_2)=p(x_1-x_2),\\
 \text{两式相加：}\quad& y_0y_m=p(x_0+x_m).
\end{aligned}
\]
当 $y_0\ne0$ 时，\textbf{相减确定方向} $k_{AB}=p/y_0$；\textbf{相加建立中点关系}，再由点斜式还原切点弦方程。当 $y_0=0$ 时，弦是竖直线 $x=-x_0$，不可强行使用斜率公式。

\end{statementbox}
